% section: 指数対数

\section{指数と対数}\label{節：指数対数}
  積・商・累乗を和と差や定数倍に変える法則を確認する.\sectionref*{節：対数型}で使う.

  \subsection{指数}\label{節：指数}
    $a$を正の整数$p$回掛けたものを$a^p$と書き,$p$を指数という.
    $a>0$なら$a^0=1$, $a^{-p}=\dfrac{1}{a^p}$と定め,有理数・実数の指数へ拡張できる.
    $f(x)=a^x$を指数関数という.実数の指数を扱うときは$a>0$とする.

  \subsection{対数}\label{節：対数}
    $a>0$, $a\ne1$, $q>0$のとき
    \[
      a^p=q\iff\log_a q=p.
    \]
    $\log$をログと読み,$p$を\bookterm{logarithm}{対数},$a$を底,$q$を真数という.
    ネイピア数
    \[
      e=\lim_{x\to\infty}\left(1+\frac1x\right)^x
    \]
    を底とするものを自然対数という.
    底$a,c$は正で$1$ではなく,真数$p,M,N,b$は正とする.
    \begin{align*}
      \log_a(p^q)&=q\log_a p,\\
      \log_a(MN)&=\log_a M+\log_a N,\\
      \log_a\left(\dfrac{M}{N}\right)&=\log_a M-\log_a N,\\
      \log_a b&=\frac{\log_c b}{\log_c a}.
    \end{align*}

    $f(x)=\log_a x$は指数関数の\bookterm{inverse}{逆関数}である.
    逆関数は,出力から元の入力を一つに戻す関数である.
    $f$の入力の範囲を\textbf{定義域}$D$,取る値の範囲を\textbf{値域}$E$とする.
    各$y\in E$について$y=f(x)$を満たす$x\in D$がただ一つなら,その$x$を返す関数$f^{-1}$を作れる.\footnote{$f^{-1}(y)$は逆関数の値であり,$\dfrac{1}{f(y)}$ではない.}
    その定義域は$E$,値域は$D$となる.点$(x,y)$が$(y,x)$へ移るので,グラフは$y=x$について対称である.

    $x\mapsto a^x$では,正の値$q$から指数$x$を一つに戻せる.その逆が$q\mapsto\log_a q$である.
    一方,$x^2$では$4$から$\pm2$へ戻れてしまう.$x\ge0$に限ると逆関数は$\sqrt y$になる.
    逆関数には,式だけでなく入力の範囲も必要なのである.
